\section{Related Work}
\label{related_work}

\subsection{Executable Requirements}
\label{sec:rw-executable}

Executable requirements aim to reduce the gap between human-readable requirement statements and machine-actionable behavior.
Controlled natural language and requirement patterns improve precision while preserving readability~\cite{darifControlledNaturalLanguage2026}.
However, such specifications do not always provide task sequences or execution dependencies for intelligent agents.

\subsection{Requirement-to-Task Transformation}
\label{sec:rw-req-task}

Requirement-to-task transformation studies how requirement statements can be converted into concrete development, validation, or operational tasks.
For this paper, the relevant concern is the production of task-level actions that preserve the original requirement intent.

\subsection{Task Planning for Actionable AI Systems}
\label{sec:rw-task-planning}

Task planning provides mechanisms for decomposing goals into ordered actions and checking whether a plan can satisfy a target objective.
In intelligent software systems, planning is useful when a requirement must be operationalized as a workflow rather than kept as a static textual artifact.

\subsection{LLM-based Requirement Engineering}
\label{sec:rw-llm-re}

LLM-based methods have been used to extract, refine, and generate software requirements.
Requirement-aware generation has shown that explicitly modeling requirements can improve downstream software artifacts~\cite{hanArchCodeIncorporatingSoftware2024}.
This paper focuses on a narrower question: how generated or stakeholder-written requirements can be transformed into executable task workflows.

\subsection{Multi-Agent Collaboration}
\label{sec:rw-multi-agent}

Multi-agent collaboration can separate analysis responsibilities across specialized roles.
For executable requirements, this separation supports independent semantic analysis, planning, and execution-flow inspection.
The proposed work uses multi-agent collaboration as a task-planning mechanism, not as a general-purpose platform architecture.
